\documentclass{article}
\usepackage[utf8]{inputenc}
\usepackage{amsmath}
\usepackage{graphicx}
\usepackage{hyperref}
\usepackage{booktabs}
\usepackage{geometry}
\usepackage[numbers]{natbib}
\geometry{margin=1in}

\title{\textbf{Polyinverse}: Multi-task Attentive Fingerprint Graph Neural Network for Polymer Property Prediction and SA Score-Guided Inverse Design}

\author{
  Ethan Im \\
  Independent Researcher, South Korea \\
  \texttt{lki100500@gmail.com}
}

\date{\today}

\begin{document}

\maketitle

\begin{abstract}
Polymer discovery traditionally relies on time-consuming trial-and-error experimentation.
We present \textbf{Polyinverse}, an open-source framework combining multi-task Attentive Fingerprint (AttentiveFP) Graph Neural Networks for polymer property prediction with a synthesizability-guided inverse design pipeline. By combining publicly available datasets augmented with pseudo-labeled structures from PI1M into a unified corpus of 7,588 unique polymer structures, our multi-task GNN simultaneously predicts three key properties:
density ($R^2 = 0.871$, MAE $= 0.0336$ g/ml), normalized crystallization temperature Tc
($R^2 = 0.728$, MAE $= 0.0262$ in normalized scale), and glass transition temperature Tg
($R^2 = 0.793$, MAE $= 68.43$ °C) from SMILES representations.
Multi-task learning improves density prediction by 35\% over single-task GNN baselines
($R^2$: $0.595 \rightarrow 0.871$), supported by a statistically significant negative
correlation between density and Tc ($r = -0.246$, $p = 1.07 \times 10^{-20}$).
For inverse design, we propose a genetic algorithm incorporating SA Score filtering
($\leq 4.0$) to ensure synthetic accessibility of candidate polymers, achieving target
density errors below 0.01 g/ml. All code and models are publicly available.
\end{abstract}

\section{Introduction}

Polymeric materials are fundamental to modern industry, appearing in applications
ranging from packaging and electronics to biomedical devices. However, developing
new polymers with desired properties remains slow and expensive, dominated by
expert intuition and laboratory trial-and-error. A single design-synthesize-test
cycle can take months, severely limiting the pace of materials innovation.

Machine learning (ML) has emerged as a powerful tool to accelerate materials discovery.
Graph neural networks (GNNs) have shown strong performance in predicting molecular
properties by directly learning from molecular graph representations,
eliminating the need for hand-crafted features \citep{gcnn_polymer_2018}.

Despite progress in property prediction, most existing work focuses on the
\textit{forward} problem: given a molecular structure, predict its properties.
The \textit{inverse} problem---given desired properties, suggest candidate
structures---remains significantly more challenging for polymers.
Existing inverse design approaches either require expensive molecular dynamics
(MD) simulations for training data \citep{vitrimer_inverse_2023} or produce
molecules with poor synthetic accessibility.

In this work, we present \textbf{Polyinverse}, addressing both problems:
\begin{enumerate}
    \item A multi-task AttentiveFP GNN that simultaneously predicts density, crystallization
    temperature (Tc), and glass transition temperature (Tg) from polymer SMILES,
    leveraging their physical correlations across a combined dataset of 7,588 polymers.
    \item A genetic algorithm-based inverse design pipeline with SA Score
    \citep{ertl_sa_score_2009} filtering for synthetic accessibility.
\end{enumerate}

\noindent\textbf{Why Genetic Algorithm over VAE/Diffusion?}
Recent generative approaches (VAE, Diffusion) require large labeled datasets
(often $>10^5$ samples) for stable training. Given the inherent scarcity of
labeled polymer data (5,616 density samples in our combined dataset), we chose a
genetic algorithm that operates directly on existing polymer SMILES without
requiring a separate generative model. This data-efficient approach is more
appropriate for the small-data regime common in polymer informatics.

\noindent Our key contributions are:
\begin{itemize}
    \item Multi-task AttentiveFP GNN achieving $R^2=0.871$ for polymer density
    using a dataset of 7,588 unique polymers augmented via PI1M pseudo-labeling
    \item Extension to three-property prediction (Density + Tc + Tg) with
    masked multi-task loss handling missing property values
    \item Statistical analysis showing Density-Tc correlation ($r=-0.246$)
    explains multi-task learning gains
    \item First integration of SA Score constraints in polymer inverse design
    via genetic algorithm, with explicit discussion of limitations for
    polymeric repeat units
    \item Fully open-source implementation with interactive web deployment
\end{itemize}

\section{Related Work}

\subsection{Polymer Property Prediction}

Early ML approaches relied on hand-crafted molecular descriptors with traditional
ML models. Reyes-Lillo et al. \citep{gcnn_polymer_2018} pioneered graph convolutional neural networks
for polymer property prediction, demonstrating superior performance over
descriptor-based methods for dielectric properties.

Multi-task learning has been explored in polymer informatics.
POLYMERGNN \citep{polymergnn_2023} demonstrated multi-task GNN prediction
of thermal and mechanical properties, achieving $R^2 \approx 0.72$ for glass
transition temperature on 240 synthesized polyesters.
The NeurIPS 2025 Open Polymer Challenge \citep{neurips_opc_2025} established
the first large-scale open benchmark, revealing that GNN-based approaches
consistently outperform descriptor-based methods.

\subsection{Polymer Inverse Design}

Yue et al. \citep{vitrimer_inverse_2023} proposed a graph VAE for inverse design of vitrimers
with target $T_g$, requiring one million MD-simulated structures for training.
Vogel and Weber \citep{copolymer_inverse_2024} extended inverse design to complex copolymer
architectures using generative models.
These approaches share a common limitation: large computational requirements
and no explicit synthesizability constraints. Our work addresses both by using
only open experimental data and enforcing SA Score filtering.

\section{Methods}

\subsection{Dataset}

We construct a unified polymer dataset combining three publicly available sources:
(1) the NeurIPS 2025 Open Polymer Challenge dataset \citep{neurips_opc_2025} containing
3,502 polymer structures, (2) the Extended Polymer Dataset with 1,088 structures,
and (3) the TG-Density Excerpt with 194 structures. After canonicalizing SMILES
representations using RDKit and removing duplicates, the experimental corpus contains
4,594 unique polymer structures.

To address the severe data scarcity in Tg (only 368 experimental samples),
we augment the dataset using pseudo-labeling from the PI1M database \citep{pi1m_2020}.
Specifically, we sample 3,000 structurally valid polymers (SA Score $\leq 4.0$)
from PI1M and generate pseudo-labels using our pre-trained model, expanding
Tg coverage from 368 to 3,362 samples. The final corpus contains 7,588 unique
polymer structures. Table~\ref{tab:dataset} summarizes the data distribution.

\begin{table}[h]
\centering
\caption{Combined dataset statistics (after PI1M augmentation)}
\label{tab:dataset}
\begin{tabular}{lcccc}
\toprule
Property & Samples & Min & Max & Mean \\
\midrule
Density (g/ml) & 5,616 & 0.135 & 1.914 & 1.094 \\
Tc (normalized) & 5,446 & 0.024 & 0.954 & 0.254 \\
Tg (°C)        & 3,362 & -138.8 & 484.7 & 156.0 \\
\bottomrule
\end{tabular}
\end{table}

\noindent\textbf{Note on Tc values:} The Tc values are provided in a normalized
scale $[0, 1]$. All reported Tc metrics are in this normalized scale.
Denormalized MAE is 0.0028 (original unit).

Polymer structures are represented as SMILES strings with repeat unit notation,
where \texttt{*} symbols at chain termini indicate polymerization attachment points.

\subsection{Multi-task Graph Neural Network}

\noindent\textbf{Molecular Representation.}
Each polymer repeat unit is represented as a molecular graph $G = (V, E)$,
where nodes $V$ correspond to atoms and edges $E$ to bonds.
Each atom $v_i$ is represented by an 8-dimensional feature vector:

\begin{equation}
\mathbf{x}_i = [\frac{Z_i}{100}, \frac{d_i}{10}, \frac{q_i}{5},
a_i, r_i, \frac{h_i}{8}, \frac{m_i}{200}, \frac{\text{hyb}_i}{8}]
\end{equation}

where $Z_i$ is atomic number, $d_i$ is degree, $q_i$ is formal charge,
$a_i \in \{0,1\}$ is aromaticity, $r_i \in \{0,1\}$ is ring membership,
$h_i$ is number of hydrogens, $m_i$ is atomic mass, and
$\text{hyb}_i$ is hybridization state (encoded as RDKit integer: SP=2, SP2=3, SP3=4).
Each bond $e_{ij}$ is represented by a 3-dimensional edge feature vector
encoding bond type, aromaticity, and ring membership.

\noindent\textbf{Architecture.}
We employ an Attentive Fingerprint (AttentiveFP) encoder \citep{attentivefp_2019}
with 4 message-passing layers and 2 readout timesteps, which applies graph-level
attention to capture both local and global molecular interactions:

\begin{equation}
\mathbf{h} = \text{AttentiveFP}(\mathbf{X}, \mathbf{E}, \mathbf{A})
\end{equation}

The shared molecular embedding passes through a fully-connected layer before
task-specific prediction heads for all three properties:

\begin{equation}
\hat{y}_{\text{density}}, \hat{y}_{T_c}, \hat{y}_{T_g} = 
f_{\text{density}}(\mathbf{h}),\; f_{T_c}(\mathbf{h}),\; f_{T_g}(\mathbf{h})
\end{equation}

\noindent\textbf{Training.}
All target properties are z-score normalized before training.
Training uses masked MSE loss to handle missing property values:

\begin{equation}
\mathcal{L} = \frac{1}{|M_d|}\sum_{i \in M_d}(\hat{y}^i_d - y^i_d)^2 +
              \frac{1}{|M_t|}\sum_{i \in M_t}(\hat{y}^i_t - y^i_t)^2 +
              \frac{1}{|M_g|}\sum_{i \in M_g}(\hat{y}^i_g - y^i_g)^2
\end{equation}

where $M_d$, $M_t$, and $M_g$ are the index sets of samples with observed
density, Tc, and Tg values, respectively.
The model is trained for 500 epochs using Adam optimizer
(lr=$10^{-3}$, weight decay=$10^{-5}$) with ReduceLROnPlateau scheduling.

\subsection{SA Score-Guided Inverse Design}

\noindent\textbf{Motivation.}
The SA Score \citep{ertl_sa_score_2009} was originally designed for
drug-like small molecules. Applying it to polymer repeat units (with \texttt{*}
terminal atoms) requires careful consideration: the \texttt{*} atoms are
treated as carbon atoms (atomic number 6) during SA Score computation in
RDKit's implementation, which may slightly underestimate synthesis difficulty
for certain repeat units. Despite this limitation, SA Score remains the most
widely accessible synthesizability metric, and we use it as a first-order
filter acknowledging that full polymerizability assessment (e.g., checking
for reactive end groups capable of chain-growth or step-growth polymerization)
is beyond the scope of this initial framework.

\noindent\textbf{Algorithm.}
Given target properties $\rho^*$, $T_c^*$, $T_g^*$, we evolve a population
of polymer SMILES using a combined fitness function:

\begin{equation}
\text{fitness}(s) = w_d|\hat{\rho}(s) - \rho^*| + w_t|\hat{T}_c(s) - T_c^*| +
w_g|\hat{T}_g(s) - T_g^*| + \lambda \cdot \max(0, \text{SA}(s) - \tau)
\end{equation}

where $w_d, w_t, w_g$ are user-defined property weights,
$\tau = 4.0$ is the synthesizability threshold, and $\lambda = 0.1$ is a penalty weight.

Mutation operators are restricted to atomic substitution within $\{C, N, O\}$
(atomic numbers 6, 7, 8) to maintain chemical reasonableness and avoid
introducing elements (F, Cl, S) that could artificially inflate density
predictions without reflecting synthetically practical polymers.
Only candidates satisfying $\text{SA}(s) \leq 4.0$ are accepted as final results.

\begin{figure}[h]
\centering
\includegraphics[width=\textwidth]{paper_figures/fig1_data_analysis.png}
\caption{Dataset analysis. (Left) Density distribution (n=5,616, mean=1.094 g/ml). 
(Center) Tc distribution in normalized scale (n=5,446). 
(Right) Tg distribution (n=3,362, after PI1M augmentation).}
\label{fig:data}
\end{figure}

\section{Results}

\subsection{Why Multi-task Learning Works: Density-Tc Correlation}

Before presenting prediction results, we analyze the physical basis for
multi-task learning gains. Figure~\ref{fig:data} shows the relationship
between density and Tc across polymers with both measurements.
The Pearson correlation coefficient $r = -0.246$ ($p = 1.07 \times 10^{-20}$)
indicates a statistically significant negative correlation: denser polymers
tend to have lower crystallization temperatures. This physical relationship
provides complementary information that multi-task learning can exploit,
explaining the substantial performance improvement over single-task models.

\subsection{Property Prediction}

Table~\ref{tab:results} compares our multi-task AttentiveFP GNN against baseline models.
All models use an 80/20 train-test split with random seed 42.

\begin{figure}[h]
\centering
\includegraphics[width=\textwidth]{paper_figures/combined_results.png}
\caption{Multi-task AttentiveFP GNN prediction results on test set. 
(Left) Density ($R^2=0.871$, MAE$=0.0336$ g/ml). 
(Center) Tc ($R^2=0.728$, MAE$=0.0262$ normalized).
(Right) Tg ($R^2=0.793$, MAE$=68.43$ °C).}
\label{fig:prediction}
\end{figure}

\begin{table}[h]
\centering
\caption{Polymer property prediction performance comparison}
\label{tab:results}
\begin{tabular}{lccccc}
\toprule
Model & Density $R^2$ & Density MAE & Tc $R^2$ & Tc MAE & Tg $R^2$ \\
\midrule
Random Forest (9 feat.)  & 0.388 & 0.0622 & --    & --     & -- \\
XGBoost (217 feat.)      & 0.595 & 0.0402 & --    & --     & -- \\
Single-task GNN          & 0.595 & 0.0364 & --    & --     & -- \\
Multi-task GNN (NeurIPS only) & 0.757 & 0.0294 & 0.512 & 0.0364 & 0.321 \\
\textbf{AttentiveFP Multi-task (ours)} & \textbf{0.871} & \textbf{0.0336} & \textbf{0.728} & \textbf{0.0262} & \textbf{0.793} \\
\bottomrule
\end{tabular}
\end{table}

Multi-task learning with the combined dataset provides a 35\% improvement
in $R^2$ over single-task GNN ($0.595 \rightarrow 0.871$).
The Tg $R^2$ improves dramatically from 0.321 to 0.793 following PI1M-based
data augmentation (368 $\rightarrow$ 3,362 samples), confirming that
data scarcity was the primary bottleneck for Tg prediction.

\subsection{Inverse Design with SA Score Filter}

Table~\ref{tab:inverse} presents representative inverse design results for
three target densities. All reported candidates satisfy SA Score $\leq 4.0$.

\begin{table}[h]
\centering
\caption{Inverse design results with SA Score filter ($\leq 4.0$)}
\label{tab:inverse}
\begin{tabular}{lcccc}
\toprule
Target (g/ml) & Description & Predicted (g/ml) & Error & SA Score \\
\midrule
0.95 & PP/PE range & 0.9532 & 0.0032 & 3.74 \\
1.20 & PET range   & 1.2099 & 0.0099 & 3.51 \\
1.40 & PVC range   & 1.3893 & 0.0107 & 3.92 \\
\bottomrule
\end{tabular}
\end{table}

Without the SA Score filter, candidates contained excessive halogen and
sulfur atoms (SA Score $> 5.5$) that are chemically unrealistic for
practical polymer synthesis. The filter effectively restricts the search
to carbon-, nitrogen-, and oxygen-based structures consistent with
commercially relevant polymer classes.

\section{Conclusion}

We presented Polyinverse, a framework for polymer property prediction
and inverse design using publicly available data.
Our multi-task AttentiveFP GNN achieves $R^2 = 0.871$ for density prediction,
outperforming single-task baselines by 35\%.
The model simultaneously predicts Tc ($R^2=0.728$) and Tg ($R^2=0.793$),
where Tg performance improved dramatically (0.321 $\rightarrow$ 0.793)
through PI1M-based pseudo-label augmentation.
The SA Score-guided genetic algorithm produces synthesizable polymer
candidates with target density errors below 0.01 g/ml, now extended
to simultaneous multi-property optimization (Density + Tc + Tg).

\noindent\textbf{Limitations.}
Pseudo-labeled Tg data from PI1M may introduce distribution shift;
future work should validate against additional experimental measurements.
The SA Score was not designed for polymeric repeat units; future work
should incorporate polymerizability-specific metrics.

\noindent\textbf{Future Work.}
We plan to (1) validate pseudo-labeled Tg predictions against experimental
data from PolyInfo and other polymer databases,
(2) explore diffusion model-based molecular generation for larger-data regimes,
and (3) incorporate a rule-based polymerizability filter to ensure
candidate repeat units possess reactive end groups capable of
chain-growth or step-growth polymerization.

The inverse design targets in this work (PP/PE range: 0.95 g/ml, 
PET range: 1.20 g/ml, PVC range: 1.40 g/ml) directly correspond to 
commercially relevant polymer classes, demonstrating that Polyinverse 
addresses real industrial materials design challenges faced by polymer 
manufacturers.

Polyinverse is publicly available at:
\url{https://github.com/Ethan-Im/polyinverse} and
\url{https://huggingface.co/spaces/Ethan-Im/polyinverse}.

\bibliographystyle{plainnat}
\begin{thebibliography}{9}

\bibitem{gcnn_polymer_2018}
Reyes-Lillo, S.E., et al. (2018).
\textit{Graph Convolutional Neural Networks for Polymer Property Prediction}.
arXiv:1811.06231.

\bibitem{polymergnn_2023}
Queen, O., et al. (2023).
\textit{POLYMERGNN: A Graph Neural Network Framework for Multi-task Polymer Property Prediction}.

\bibitem{neurips_opc_2025}
Liu, G., et al. (2025).
\textit{Open Polymer Challenge: Post-Competition Report}.
arXiv:2512.08896.

\bibitem{vitrimer_inverse_2023}
Yue, Y., et al. (2023).
\textit{AI-guided inverse design and discovery of recyclable vitrimeric polymers}.
arXiv:2312.03690.

\bibitem{copolymer_inverse_2024}
Vogel, G., Weber, J.M. (2024).
\textit{Inverse design of copolymers including stoichiometry and chain architecture}.
Chemical Science. DOI:10.1039/d4sc05900j.

\bibitem{ertl_sa_score_2009}
Ertl, P., Schuffenhauer, A. (2009).
\textit{Estimation of synthetic accessibility score of drug-like molecules}.
Journal of Cheminformatics, 1(1), 8.

\bibitem{attentivefp_2019}
Xiong, Z., et al. (2019).
\textit{Pushing the Boundaries of Molecular Representation for Drug Discovery with the Graph Attention Mechanism}.
Journal of Medicinal Chemistry, 63(16), 8749--8760.

\bibitem{pi1m_2020}
Ma, R., Luo, T. (2020).
\textit{PI1M: A Benchmark Database for Polymer Informatics}.
Journal of Chemical Information and Modeling, 61(6), 2712--2721.

\bibitem{kipf_gcn_2017}
Kipf, T.N., Welling, M. (2017).
\textit{Semi-Supervised Classification with Graph Convolutional Networks}.
ICLR 2017. arXiv:1609.02907.

\end{thebibliography}

\end{document}
